% 凸集与非凸集：(a) 椭圆盘内任意两点的线段留在集合内；(b) 月牙形中有线段穿出集合；
% (c) 资源分配可行多边形 x1+x2<=4, 2x1+x2<=6, x>=0 中两种方案 P=(3,0), Q=(0,4) 的等权混合 M=(1.5,2) 仍可行。
\begin{tikzpicture}[note/.style={font=\FigureNoteFont\fontsize{7.5}{9}\selectfont,align=center},
  pt/.style={circle,fill=#1,inner sep=0pt,minimum size=3.2pt}]
  % (a) convex
  \begin{scope}
    \draw[FigureBlue,fill=FigureBlueFill,line width=.9pt,rotate=20] (0,0) ellipse (1.25 and .8);
    \coordinate (u) at (-.85,-.45); \coordinate (v) at (.9,.55);
    \draw[FigureGreen,line width=.9pt] (u)--(v);
    \node[pt=FigureGreen] at (u) {}; \node[pt=FigureGreen] at (v) {};
    \node[below left=-1pt,text=FigureGreen] at (u) {$\bm u$};
    \node[above right=-1pt,text=FigureGreen] at (v) {$\bm v$};
    \node[note] at (0,-1.45) {(a) 凸集};
  \end{scope}
  % (b) nonconvex crescent
  \begin{scope}[xshift=3.3cm]
    \draw[FigureBlue,fill=FigureBlueFill,line width=.9pt]
      (0.4583,0.9447) arc[start angle=64.12,end angle=326.39,radius=1.05]
      arc[start angle=293.93,end angle=96.59,radius=.8] -- cycle;
    \coordinate (u) at (.15,.98); \coordinate (v) at (.2,-.95);
    \draw[FigureRed,line width=.9pt,dashed] (u)--(v);
    \node[pt=FigureRed] at (u) {}; \node[pt=FigureRed] at (v) {};
    \node[above=0pt,text=FigureRed] at (u) {$\bm u$};
    \node[below=0pt,text=FigureRed] at (v) {$\bm v$};
    \node[note] at (0,-1.45) {(b) 非凸集};
  \end{scope}
  % (c) resource polygon
  \begin{scope}[xshift=5.6cm,yshift=-1.0cm,scale=.55]
    \draw[->,FigureStroke,line width=.6pt] (-.3,0)--(3.8,0) node[right]{$x_1$};
    \draw[->,FigureStroke,line width=.6pt] (0,-.3)--(0,4.5) node[above]{$x_2$};
    \draw[FigureBlue,fill=FigureBlueFill,line width=.9pt] (0,0)--(3,0)--(2,2)--(0,4)--cycle;
    \draw[FigureGreen,line width=.9pt] (3,0)--(0,4);
    \node[pt=FigureGreen] at (3,0) {}; \node[pt=FigureGreen] at (0,4) {};
    \node[pt=FigureInk] at (1.5,2) {};
    \node[below=1pt,text=FigureGreen] at (3,0) {$P$};
    \node[left=1pt,text=FigureGreen] at (0,4) {$Q$};
    \node[left=1pt,text=FigureInk] at (1.45,1.95) {$M$};
    \node[note] at (1.6,-.95) {(c) 可行多边形};
  \end{scope}
\end{tikzpicture}
